\documentclass[12pt,letterpaper]{article}
\usepackage{graphicx,textcomp}
\usepackage{natbib}
\usepackage{setspace}
\usepackage[margin=1in]{geometry}
\usepackage{color}
\usepackage[reqno]{amsmath}
\usepackage{amsthm}
\usepackage{fancyvrb}
\usepackage{amssymb,enumerate}
\usepackage[all]{xy}
\usepackage{endnotes}
\usepackage{lscape}
\newtheorem{com}{Comment}
\usepackage{float}
\usepackage{hyperref}
\newtheorem{lem} {Lemma}
\newtheorem{prop}{Proposition}
\newtheorem{thm}{Theorem}
\newtheorem{defn}{Definition}
\newtheorem{cor}{Corollary}
\newtheorem{obs}{Observation}
\usepackage[compact]{titlesec}
\usepackage{dcolumn}
\usepackage{tikz}
\usetikzlibrary{arrows}
\usepackage{multirow}
\usepackage{xcolor}
\newcolumntype{.}{D{.}{.}{-1}}
\newcolumntype{d}[1]{D{.}{.}{#1}}
\definecolor{light-gray}{gray}{0.65}
\usepackage{url}
\usepackage{listings}
\usepackage{color}

\definecolor{codegreen}{rgb}{0,0.6,0}
\definecolor{codegray}{rgb}{0.5,0.5,0.5}
\definecolor{codepurple}{rgb}{0.58,0,0.82}
\definecolor{backcolour}{rgb}{0.95,0.95,0.92}

\lstdefinestyle{mystyle}{
	backgroundcolor=\color{backcolour},   
	commentstyle=\color{codegreen},
	keywordstyle=\color{magenta},
	numberstyle=\tiny\color{codegray},
	stringstyle=\color{codepurple},
	basicstyle=\footnotesize,
	breakatwhitespace=false,         
	breaklines=true,                 
	captionpos=b,                    
	keepspaces=true,                 
	numbers=left,                    
	numbersep=5pt,                  
	showspaces=false,                
	showstringspaces=false,
	showtabs=false,                  
	tabsize=2
}
\lstset{style=mystyle}
\newcommand{\Sref}[1]{Section~\ref{#1}}
\newtheorem{hyp}{Hypothesis}

\title{Problem Set 3}
\date{Due: March 25, 2026}
\author{Applied Stats II}


\begin{document}
	\maketitle
	\section*{Instructions}
	\begin{itemize}
	\item Please show your work! You may lose points by simply writing in the answer. If the problem requires you to execute commands in \texttt{R}, please include the code you used to get your answers. Please also include the \texttt{.R} file that contains your code. If you are not sure if work needs to be shown for a particular problem, please ask.
\item Your homework should be submitted electronically on GitHub in \texttt{.pdf} form.
\item This problem set is due before 23:59 on Wednesday March 25, 2026. No late assignments will be accepted.
	\end{itemize}

	\vspace{.25cm}
\section*{Question 1}
\vspace{.25cm}
\noindent We are interested in how governments' management of public resources impacts economic prosperity. Our data come from \href{https://www.researchgate.net/profile/Adam_Przeworski/publication/240357392_Classifying_Political_Regimes/links/0deec532194849aefa000000/Classifying-Political-Regimes.pdf}{Alvarez, Cheibub, Limongi, and Przeworski (1996)} and is labelled \texttt{gdpChange.csv} on GitHub. The dataset covers 135 countries observed between 1950 or the year of independence or the first year forwhich data on economic growth are available ("entry year"), and 1990 or the last year for which data on economic growth are available ("exit year"). The unit of analysis is a particular country during a particular year, for a total $>$ 3,500 observations. 

\begin{itemize}
	\item
	Response variable: 
	\begin{itemize}
		\item \texttt{GDPWdiff}: Difference in GDP between year $t$ and $t-1$. Possible categories include: "positive", "negative", or "no change"
	\end{itemize}
	\item
	Explanatory variables: 
	\begin{itemize}
		\item
		\texttt{REG}: 1=Democracy; 0=Non-Democracy
		\item
		\texttt{OIL}: 1=if the average ratio of fuel exports to total exports in 1984-86 exceeded 50\%; 0= otherwise
	\end{itemize}
	
\end{itemize}
\newpage
\noindent Please answer the following questions:

\begin{enumerate}
	\item Construct and interpret an unordered multinomial logit with \texttt{GDPWdiff} as the output and "no change" as the reference category, including the estimated cutoff points and coefficients.
	\item Construct and interpret an ordered multinomial logit with \texttt{GDPWdiff} as the outcome variable, including the estimated cutoff points and coefficients.
	
	
\end{enumerate}

\section*{Question 2} 
\vspace{.25cm}

\noindent Consider the data set \texttt{MexicoMuniData.csv}, which includes municipal-level information from Mexico. The outcome of interest is the number of times the winning PAN presidential candidate in 2006 (\texttt{PAN.visits.06}) visited a district leading up to the 2009 federal elections, which is a count. Our main predictor of interest is whether the district was highly contested, or whether it was not (the PAN or their opponents have electoral security) in the previous federal elections during 2000 (\texttt{competitive.district}), which is binary (1=close/swing district, 0="safe seat"). We also include \texttt{marginality.06} (a measure of poverty) and \texttt{PAN.governor.06} (a dummy for whether the state has a PAN-affiliated governor) as additional control variables. 

\begin{enumerate}
	\item [(a)]
	Run a Poisson regression because the outcome is a count variable. Is there evidence that PAN presidential candidates visit swing districts more? Provide a test statistic and p-value.

	\item [(b)]
	Interpret the \texttt{marginality.06} and \texttt{PAN.governor.06} coefficients.
	
	\item [(c)]
	Provide the estimated mean number of visits from the winning PAN presidential candidate for a hypothetical district that was competitive (\texttt{competitive.district}=1), had an average poverty level (\texttt{marginality.06} = 0), and a PAN governor (\texttt{PAN.governor.06}=1).
	
\end{enumerate}
\newpage
\section*{Solution - Question 1}
\vspace{.25cm}
Before fitting the models, the outcome was transformed into a categorical variable. The values equal to 0 were assigned to the \texttt{no\_change} category, everything below that to the \texttt{negative} category, and everything above to the \texttt{positive} category. Then, \texttt{no\_change} was selected as the reference category, as it makes the most intuitive sense in the interpretation. 
\vspace{0.25cm}
\lstinputlisting[language=R, firstline=39,lastline=47]{PS3_DP.R}
\begin{enumerate}
	\item \textbf{Unordered multinomial logit}
	\vspace{.25cm}
		\lstinputlisting[language=R, firstline=49,lastline=57]{PS3_DP.R}
		\begin{table}[!htbp] \centering 
			\caption{Regression Summary 1} 
			\label{} 
			\begin{tabular}{@{\extracolsep{5pt}}lcc} 
				\\[-1.8ex]\hline 
				\hline \\[-1.8ex] 
				& \multicolumn{2}{c}{\textit{Dependent variable:}} \\ 
				\cline{2-3} 
				\\[-1.8ex] & negative change in GDP & positive change in GDP \\ 
				\\[-1.8ex] & (1) & (2)\\ 
				\hline \\[-1.8ex] 
				democracy & 1.379$^{*}$ & 1.769$^{**}$ \\ 
				& (0.769) & (0.767) \\ 
				& & \\ 
				fuel exports ratio above 50\% & 4.784 & 4.576 \\ 
				& (6.885) & (6.885) \\ 
				& & \\ 
				constant & 3.805$^{***}$ & 4.534$^{***}$ \\ 
				& (0.271) & (0.269) \\ 
				& & \\ 
				\hline \\[-1.8ex] 
				\textit{Note:}  & \multicolumn{2}{r}{$^{*}$p$<$0.1; $^{**}$p$<$0.05; $^{***}$p$<$0.01} \\ 
			\end{tabular} 
		\end{table} 
		\newpage
		\begin{table}[!htbp] \centering 
			\caption{Exponentiated Coefficients 1} 
			\label{} 
			\begin{tabular}{@{\extracolsep{5pt}} cccc} 
				\\[-1.8ex]\hline 
				\hline \\[-1.8ex] 
				& (Intercept) & democracy & fuel exports ratio above 50\% \\ 
				\hline \\[-1.8ex] 
				negative change in GDP & $44.942$ & $3.972$ & $119.578$ \\ 
				positive change in GDP & $93.108$ & $5.865$ & $97.156$ \\ 
				\hline \\[-1.8ex] 
			\end{tabular} 
		\end{table} 
		\textbf{\underline{Interpretation }}
		\begin{itemize}
			\item \textbf{Intercept:} \\[0.2cm] The baseline log-odds (or the log-odds of a country that is not a democracy with a ratio of fuel exports to total exports that is less than 50\%) of observing a negative change in GDP, compared to observing no change in GDP between year t and t-1 is 3.805. Similarly, the baseline log-odds (non-democracy, fuel exports ratio less than 50\%) of observing a positive change in GDP, compared to observing no change in GDP between year t and year t-1 is 4.534. The baseline odds of observing a negative GDP change relative to no change are approximately 44.9, while the baseline odds of observing a positive GDP change are approximately 93.1.
			\item \textbf{Democracy:} \\[0.2cm] Holding all other variables constant, moving from a non-democracy to a democracy increases the log-odds of observing a negative change in GDP, compared to no change in GDP, by 1.379, and the log-odds of observing a positive change in GDP, compared to no change in GDP, by 1.769. In terms of odds, holding all other variables constant, moving from a non-democracy to a democracy increases the odds of observing a negative change in GDP (relative to no change) by a factor of 3.97 and the odds of observing a positive change by a factor of 5.87. Both effects are statistically significant, but the negative change in GDP coefficient only at the 10\% level, while the positive change is also significant at the 5\% level, suggesting a more reliable relationship.
			\item \textbf{Fuel exports ratio above 50\%: } \\[0.2cm] Holding regime constant, moving from a ratio of fuel exports to total exports that is less than 50\% to one that is more than 50\%, increases the log-odds of observing a negative change in GDP, compared to no change in GDP, by 4.784 and the log-odds of observing a positive change in GDP, compared to no change in GDP, by 4.576. For oil dependence, holding regime constant, moving from a fuel export ratio below 50\% to above 50\% increases the odds of a negative GDP change by a factor of 119.6 and a positive GDP change by a factor of 97.2, both compared to no change in GDP. However, these effects are not statistically significant and have very large standard errors, suggesting that the estimated relationships are not reliable.
		\end{itemize}
\newpage
	\item \textbf{Ordered multinomial logit} \\[.25cm]
				Prior to estimating the ordered multinomial logit model, the outcome was relevelled and converted into an ordered factor.
		\vspace{.25cm}
		\lstinputlisting[language=R, firstline=59,lastline=75]{PS3_DP.R}
			\begin{table}[!htbp] \centering 
				\caption{Regression Summary 2} 
				\label{} 
				\begin{tabular}{@{\extracolsep{5pt}}lc} 
					\\[-1.8ex]\hline 
					\hline \\[-1.8ex] 
					& \multicolumn{1}{c}{\textit{Dependent variable:}} \\ 
					\cline{2-2} 
					\\[-1.8ex] & change in GDP \\ 
					\hline \\[-1.8ex] 
					democracy & 0.398$^{***}$ \\ 
					& (0.075) \\ 
					& \\ 
					fuel exports ratio above 50\% & $-$0.199$^{*}$ \\ 
					& (0.116) \\ 
					& \\ 
					\hline \\[-1.8ex] 
					Observations & 3,721 \\ 
					\hline 
					\hline \\[-1.8ex] 
					\textit{Note:}  & \multicolumn{1}{r}{$^{*}$p$<$0.1; $^{**}$p$<$0.05; $^{***}$p$<$0.01} \\ 
				\end{tabular} 
			\end{table} 
			\begin{table}[!htbp] \centering 
				\caption{Exponentiated Coefficients 2} 
				\label{} 
				\begin{tabular}{@{\extracolsep{5pt}} cc} 
					\\[-1.8ex]\hline 
					\hline \\[-1.8ex] 
					democracy & fuel exports ratio above 50\% \\ 
					\hline \\[-1.8ex] 
					$1.490$ & $0.820$ \\ 
					\hline \\[-1.8ex] 
				\end{tabular} 
			\end{table} 
	\newpage
	\textbf{\underline{Interpretation}}
	\begin{itemize}
		\item \textbf{Democracy: } \\[0.2cm]
		Holding all other variables constant, for democracies, the odds of observing a more positive difference (moving from negative to no change, or no change to positive) in GDP between year t and t-1 are approximately 1.49 times higher compared to non-democracies. Being a democracy is associated with a 0.398 increase in the log-odds of being in a higher (more positive) category of GDP change. The regime variable is statistically significant. 
		\item \textbf{Fuel exports ratio above 50\%:} \\[0.2cm] Holding all other variables constant, countries with an average ratio of fuel exports to total exports exceeding 50\% in the period under study have approximately 0.82 times lower odds of being in a more positive category of GDP change, compared to countries with a smaller fuel export ratio. Being a heavy fuel exporter is associated with a approximately 0.2 decrease in the log-odds of being in a higher category of GDP change. Nevertheless, this variable is only statistically significant at the 10\% level. 
			\begin{table}[!htbp] \centering 
			\caption{Cutoff Points} 
			\label{} 
			\begin{tabular}{@{\extracolsep{5pt}} cc} 
				\\[-1.8ex]\hline 
				\hline \\[-1.8ex] 
				negative\textbar no\_change & no\_change\textbar positive \\ 
				\hline \\[-1.8ex] 
				$$-$0.731$ & $$-$0.710$ \\ 
				\hline \\[-1.8ex] 
			\end{tabular} 
		\end{table} 
		\item \textbf{Cutoff points: } \\[0.2cm]
		The cutoff points are the thresholds between GDP difference categories, on an imaginary latent scale. The thresholds between the negative change and no change categories (-0.731) and the no change to positive categories (-0.710) are relatively close together, suggesting that the no change category is quite narrow; there is not a lot of space for variation for the countries that fall there. Instead, countries lean almost always towards either the negative difference or the positive difference categories. This makes sense given the coding of the data, with the no change category corresponding only to a value of 0, not allowing for a range of values close to 0. Even if it is a small difference of, for example, 0.1, due to the coding process, countries will be assigned to the positive category, which results in an inherently small range for the no change category, a range in which not a lot of countries will naturally fall.  
	\end{itemize}
\end{enumerate}
\newpage 
\section*{Solution - Question 2}
\vspace{0.25cm}
\begin{enumerate}
	\item [(a)] \textbf{Poisson regression}
		\lstinputlisting[language=R, firstline=86,lastline=97]{PS3_DP.R}
		\begin{verbatim}
			test statistic: -0.4766106 | p-value: 0.6336394 
		\end{verbatim}
		Looking at the test statistic and at the p-value for the competitive district variable, we can conclude that there is no evidence that the PAN presidential candidates visit swing districts more frequently. We can interpret this result within the framework of a hypothesis test, where the null hypothesis states that the coefficient of the competitive district variable is equal to zero. SInce the p-value (0.63) exceeds the 0.05 threshold, we fail to reject the null hypothesis. There is no evidence that the coefficient for the competitive district dummy is statistically different from zero. Thus, there is no evidence that competitive (swing) districts receive more visits from the PAN presidential candidates, compared to "safe seat" districts. 
	\item [(b)] \textbf{Interpretation}
	\begin{itemize}
		\item \texttt{marginality.06} \textbf{- measure of poverty :} \\[.2cm] Holding all other variables constant, a one-unit increase in the marginality (poverty) measure is associated with a multiplicative decrease in the expected number of visits from the winning PAN presidential candidate leading up to the federal elections by a factor of 0.12, which corresponds to a 88\% decrease in visits. This effect is statistically significant, indicating that poorer districts received, on average, fewer visits from the presidential candidate. 
		\item \texttt{PAN.governor.06} \textbf{- dummy for PAN-affiliated governor: } \\[.2cm]Holding all other variables constant, being in a district with a PAN-affiliated governor is associated with a multiplicative decrease in the expected number of visits from the PAN presidential candidate by a factor of 0.73, corresponding to a 27\% decrease in the number of visits. However, this effect is only statistically significant at the 10\% level, and not at the conventional 5\% threshold. 
	\end{itemize}
	\newpage
	\begin{table}[!htbp] \centering 
		\caption{Poisson Regression Summary} 
		\label{} 
		\begin{tabular}{@{\extracolsep{5pt}}lc} 
			\\[-1.8ex]\hline 
			\hline \\[-1.8ex] 
			& \multicolumn{1}{c}{\textit{Dependent variable:}} \\ 
			\cline{2-2} 
			\\[-1.8ex] & visits from PAN candidates \\ 
			\hline \\[-1.8ex] 
			competitive district & $-$0.081 \\ 
			& (0.171) \\ 
			& \\ 
			marginality & $-$2.080$^{***}$ \\ 
			& (0.117) \\ 
			& \\ 
			PAN-affiliated governor & $-$0.312$^{*}$ \\ 
			& (0.167) \\ 
			& \\ 
			constant & $-$3.810$^{***}$ \\ 
			& (0.222) \\ 
			& \\ 
			\hline \\[-1.8ex] 
			Observations & 2,407 \\ 
			Log Likelihood & $-$645.606 \\ 
			Akaike Inf. Crit. & 1,299.213 \\ 
			\hline 
			\hline \\[-1.8ex] 
			\textit{Note:}  & \multicolumn{1}{r}{$^{*}$p$<$0.1; $^{**}$p$<$0.05; $^{***}$p$<$0.01} \\ 
		\end{tabular} 
	\end{table} 
	\lstinputlisting[language=R, firstline=100,lastline=102]{PS3_DP.R}
	\begin{table}[!htbp] \centering 
		\caption{Exponentiated Coefficients 3} 
		\label{} 
		\begin{tabular}{@{\extracolsep{5pt}} cccc} 
			\\[-1.8ex]\hline 
			\hline \\[-1.8ex] 
			(Intercept) & marginality& PAN-affiliated governor \\ 
			\hline \\[-1.8ex] 
			$0.022$ & $0.125$ & $0.732$ \\ 
			\hline \\[-1.8ex] 
		\end{tabular} 
	\end{table} 
	\vspace{.25cm}
  	\item [(c)] \textbf{Mean number of visits prediction} \\[0.2cm]
  	To calculate the expected mean number of visits for a district with the desired characteristics, I first created a new data frame storing these values. Then, I used the predict() function, passing the previously fitted Poisson model, along with the new data frame. The type argument was set to \texttt{response} to ensure that the function returned predictions on the scale of the outcome variable, in this case, the expected number of visits from the PAN presidential candidate. 
  	\newpage
  	\lstinputlisting[language=R, firstline=105,lastline=115]{PS3_DP.R}
  	\begin{verbatim}
  		0.01494818 
  	\end{verbatim}
  	For a competitive district, with an average poverty level and a PAN governor, the estimated mean number of visits from the winning PAN presidential candidate is approximately 0.015. Substantively, on average, a district with these characteristics would esentially receive no visits from the PAN presidential candidate, highlighting very limited campaign attention for this type of districts.
\end{enumerate} 
\end{document}
